\subsection{Smith exponents on arbitrary finite $p$-adic windows}
\label{psw-subsection}

The perturbation argument in Adriazola's reconstruction
\cite[Section~12]{adriazola-reconstruction}
admits the following strengthening.  The collision term, ramified-step
extension and valuation-ring proof below refine that argument.  This is a local quantitative
lemma.  Its application still requires the actual Gram-polynomial
identity, coefficient integrality and leading-coefficient bound.

Write $v(p)=1$ on $\mathbb C_p$ and
$\mathcal O=\{a:v(a)\ge0\}$.  Suppose
\[
 A(z)=\sum_{k=0}^{\Delta}C_k\binom zk,
 \qquad C_k\in M_n(\mathcal O),\qquad
 f(z)=\det A(z)=\ell\prod_{\rho\in R}(z-\rho)\ne0,
\]
where $n\ge1$, $\Delta\ge0$, and roots are counted with
multiplicity.  Put
\[
 R_+=\{\rho\in R:v(\rho)\ge0\},\quad s=|R_+|,\quad
 c=v(\ell)+\sum_{v(\rho)<0}v(\rho),\quad
 L=\lfloor\log_p\max(1,\Delta)\rfloor.
\]
For $Y\ge1$ distinct points $z_1,\ldots,z_Y\in\mathbb Z_p$ set
\[
 C(Z)=\max_k\sum_{j\ne k}v(z_j-z_k).
\]
The matrix $A(z_j)$ is integral; let $a_1(j)\le\cdots\le a_n(j)$
be its Smith exponents, with $+\infty$ for a zero diagonal entry.

\begin{theorem}\label{psw-window}
For every real $E>0$,
\begin{equation}\label{psw-general}
 E\sum_{j=1}^Y\#\{i:a_i(j)\ge E\}
       \le Yc+s\{E+L+C(Z)\}.
\end{equation}
If $z_j=t+dj$ for $0\le j<Y$, where
$t\in\mathbb Z_p$ and $0\ne d\in\mathbb Z_p$, then
\[
 C(Z)=(Y-1)v_p(d)+v_p((Y-1)!).
\]
There is no upper restriction on $Y$.  A second valid form of
\eqref{psw-general} replaces $s$ by $\deg f$ and $c$ by $v(\ell)$.
\end{theorem}

\begin{proof}
Finite Smith elimination works over $\mathcal O$: choose an entry
of minimum valuation, divide other entries by that pivot, and clear
its row and column by integral invertible operations.  Induction
gives a diagonal matrix.  The valuations of the determinantal ideals
identify its exponents uniquely.  No discreteness assumption is used.

First let $E$ be rational.  If integral matrices $B,B'$ differ
entrywise by terms of valuation at least $E$, choose $q$ with
$v(q)=E$.  Then
\[
 B\mathcal O^n+q\mathcal O^n=B'\mathcal O^n+q\mathcal O^n.
\]
The determinant ideal of this finitely generated full-rank module
is generated by the maximal minors of $[B\ qI_n]$.  Its valuation
is $\sum_i\min(a_i(B),E)$ by Smith reduction.  Hence
\begin{equation}\label{psw-trunc}
 \sum_i\min(a_i(B),E)=\sum_i\min(a_i(B'),E).
\end{equation}

If $z\in\mathbb Z_p$ and $v(y)=E+L$, then
\[
 \binom yk=\frac yk\prod_{i=1}^{k-1}\frac{y-i}{i}
 \quad\Longrightarrow\quad
 v\binom yk=v(y)-v_p(k)\ge E\qquad(1\le k\le\Delta).
\]
Indeed $v_p(i)\le L<v(y)$.  The polynomial Vandermonde identity,
and $\binom zk\in\mathbb Z_p$, imply
$v(A(z+y)_{ab}-A(z)_{ab})\ge E$.  For $\Delta=0$ this
conclusion is automatic.

For each fixed $z=z_j$ choose such a $y$ so that
\[
 v(z+y-\rho)=\min\{v(z-\rho),E+L\}\qquad(\rho\in R).
\]
Writing $y=qu$ with $v(q)=E+L$, each root with
$v(z-\rho)=E+L$ excludes at most one residue class for the
unit $u$.  The residue field of $\mathbb C_p$ is infinite, so a
permitted $u$ exists.  For all other roots the equality is automatic.
Consequently $f(z+y)\ne0$, and \eqref{psw-trunc} gives
\[
 E\#\{i:a_i(j)\ge E\}
 \le v(f(z_j+y))
 =c+\sum_{\rho\in R_+}\min\{v(z_j-\rho),E+L\}.
\]
Roots outside the unit disk contribute their fixed valuations
because $z_j$ and $y$ are integral.

For each $\rho\in R_+$ choose $k$ maximizing $v(z_k-\rho)$.
For $j\ne k$, the ultrametric inequality gives
$v(z_j-\rho)\le v(z_j-z_k)$: a strict reverse inequality would
force $v(z_k-\rho)=v(z_j-z_k)<v(z_j-\rho)$.
The $k$-th capped term is at most $E+L$ and the remaining terms
sum to at most $C(Z)$.  Summing proves \eqref{psw-general}.
For real $E$ pass to the limit through rational thresholds
increasing to $E$; the finite lists of exponents give the same limit.

For a progression, the collision sum at $k$ is
\[
 (Y-1)v_p(d)+v_p(k!)+v_p((Y-1-k)!)
 \le (Y-1)v_p(d)+v_p((Y-1)!),
\]
by integrality of $\binom{Y-1}k$, with equality at either endpoint.
The degree-only variant follows by applying the same closest-sample
bound to every root, including those outside the unit disk.
\end{proof}

If $p\nmid d$, $\Delta<p^h$, $r=\deg f\le n\Delta$, and
$v(\ell)+r/(p-1)\le n\Lambda_0$, the theorem yields
\[
 \frac1Y\sum_j\#\{i:a_i(j)\ge E\}
 \le n\left(\frac\Delta Y+
          \frac{\Lambda_0+\Delta(h-1)/Y}{E}\right).
\]
This recovers the source's estimate without its condition $Y\le p^h$.
The exact version retains $r$ (or $s$) and the factorial valuation,
and remains valid for $p$-divisible steps after adding the displayed
collision term.  That term is necessary: for $A(z)=z$, $d=p^a$
and $E\le a$, every sampled exponent is at least $E$.

The source proof begins by putting arbitrary $\mathbb C_p$
coefficients into a finite extension of $\mathbb Q_p$; that step
is not valid at the stated generality.  The valuation-ring argument
above repairs it without restricting the coefficients.  This repair
does not supply any of the separate arithmetic family hypotheses.
